\documentclass[../../main.tex]{subfiles}
\begin{document}

{\bf [解答]}

% TODO: xy平面上に置かれたABCDなどの概要図

まず一辺の長さ$\sqrt{2}$として考え，最後に全体を$\dfrac{a}{\sqrt{2}}$倍させてスケールさせることで答えを導く．

このとき，$A(1,0)$，$B(0,-1)$，$C(-1,0)$となるような座標を定める．すると，$AD:x+y=1$であるから，$Q(1-t,t) \ (0\le t \le 1)$と置ける．すると$P$は直線$TQ$上の点だから$P(x,t)$とおける．

% TODO: 手書きpdfの
% \begin{center}
%  \scalebox{1}{%WinTpicVersion4.32a
{\unitlength 0.1in%
\begin{picture}(20.0000,20.0000)(4.0000,-24.0000)%
% STR 2 0 3 0 Black White  
% 4 1390 1397 1390 1410 4 400 0 0
% O
\put(13.9000,-14.1000){\makebox(0,0)[rt]{O}}%
% STR 2 0 3 0 Black White  
% 4 1360 387 1360 400 4 400 0 0
% $y$
\put(13.6000,-4.0000){\makebox(0,0)[rt]{$y$}}%
% STR 2 0 3 0 Black White  
% 4 2400 1427 2400 1440 4 400 0 0
% $x$
\put(24.0000,-14.4000){\makebox(0,0)[rt]{$x$}}%
% VECTOR 2 0 3 0 Black White  
% 2 1400 2400 1400 400
% 
\special{pn 8}%
\special{pa 1400 2400}%
\special{pa 1400 400}%
\special{fp}%
\special{sh 1}%
\special{pa 1400 400}%
\special{pa 1380 467}%
\special{pa 1400 453}%
\special{pa 1420 467}%
\special{pa 1400 400}%
\special{fp}%
% VECTOR 2 0 3 0 Black White  
% 2 400 1400 2400 1400
% 
\special{pn 8}%
\special{pa 400 1400}%
\special{pa 2400 1400}%
\special{fp}%
\special{sh 1}%
\special{pa 2400 1400}%
\special{pa 2333 1380}%
\special{pa 2347 1400}%
\special{pa 2333 1420}%
\special{pa 2400 1400}%
\special{fp}%
% FUNC 2 0 3 0 Black White  
% 9 400 400 2400 2400 1400 1400 2200 1400 1400 600 1400 400 2200 2400 0 4 0 0
% x-1
\special{pn 8}%
\special{pn 8}%
\special{pa 1200 2400}%
\special{pa 1206 2394}%
\special{ip}%
\special{pa 1233 2367}%
\special{pa 1239 2361}%
\special{ip}%
\special{pa 1267 2333}%
\special{pa 1273 2327}%
\special{ip}%
\special{pa 1300 2300}%
\special{pa 1306 2294}%
\special{ip}%
\special{pa 1333 2267}%
\special{pa 1339 2261}%
\special{ip}%
\special{pa 1367 2233}%
\special{pa 1373 2227}%
\special{ip}%
\special{ip}%
\special{pa 1400 2200}%
\special{pa 2200 1400}%
\special{fp}%
\special{pn 8}%
\special{pa 2206 1394}%
\special{pa 2233 1367}%
\special{ip}%
\special{pa 2239 1361}%
\special{pa 2267 1333}%
\special{ip}%
\special{pa 2273 1327}%
\special{pa 2300 1300}%
\special{ip}%
\special{pa 2306 1294}%
\special{pa 2333 1267}%
\special{ip}%
\special{pa 2339 1261}%
\special{pa 2367 1233}%
\special{ip}%
\special{pa 2373 1227}%
\special{pa 2400 1200}%
\special{ip}%
% FUNC 2 0 3 0 Black White  
% 10 400 400 2400 2400 1400 1400 2200 1400 1400 600 1400 400 2200 2400 0 4 0 0 0 0
% -x+1
\special{pn 8}%
\special{pn 8}%
\special{pa 1200 400}%
\special{pa 1206 406}%
\special{ip}%
\special{pa 1233 433}%
\special{pa 1239 439}%
\special{ip}%
\special{pa 1267 467}%
\special{pa 1273 473}%
\special{ip}%
\special{pa 1300 500}%
\special{pa 1306 506}%
\special{ip}%
\special{pa 1333 533}%
\special{pa 1339 539}%
\special{ip}%
\special{pa 1367 567}%
\special{pa 1373 573}%
\special{ip}%
\special{ip}%
\special{pa 1400 600}%
\special{pa 2200 1400}%
\special{fp}%
\special{pn 8}%
\special{pa 2206 1406}%
\special{pa 2233 1433}%
\special{ip}%
\special{pa 2239 1439}%
\special{pa 2267 1467}%
\special{ip}%
\special{pa 2273 1473}%
\special{pa 2300 1500}%
\special{ip}%
\special{pa 2306 1506}%
\special{pa 2333 1533}%
\special{ip}%
\special{pa 2339 1539}%
\special{pa 2367 1567}%
\special{ip}%
\special{pa 2373 1573}%
\special{pa 2400 1600}%
\special{ip}%
% FUNC 2 0 3 0 Black White  
% 10 400 400 2400 2400 1400 1400 2200 1400 1400 600 600 400 1400 2400 0 4 0 0 0 0
% x+1
\special{pn 8}%
\special{pn 8}%
\special{pa 400 1600}%
\special{pa 406 1594}%
\special{ip}%
\special{pa 433 1567}%
\special{pa 439 1561}%
\special{ip}%
\special{pa 467 1533}%
\special{pa 473 1527}%
\special{ip}%
\special{pa 500 1500}%
\special{pa 506 1494}%
\special{ip}%
\special{pa 533 1467}%
\special{pa 539 1461}%
\special{ip}%
\special{pa 567 1433}%
\special{pa 573 1427}%
\special{ip}%
\special{ip}%
\special{pa 600 1400}%
\special{pa 1400 600}%
\special{fp}%
\special{pn 8}%
\special{pa 1406 594}%
\special{pa 1433 567}%
\special{ip}%
\special{pa 1439 561}%
\special{pa 1467 533}%
\special{ip}%
\special{pa 1473 527}%
\special{pa 1500 500}%
\special{ip}%
\special{pa 1506 494}%
\special{pa 1533 467}%
\special{ip}%
\special{pa 1539 461}%
\special{pa 1567 433}%
\special{ip}%
\special{pa 1573 427}%
\special{pa 1600 400}%
\special{ip}%
% FUNC 2 0 3 0 Black White  
% 10 400 400 2400 2400 1400 1400 2200 1400 1400 600 600 400 1400 2400 0 4 0 0 0 0
% -x-1
\special{pn 8}%
\special{pn 8}%
\special{pa 400 1200}%
\special{pa 406 1206}%
\special{ip}%
\special{pa 433 1233}%
\special{pa 439 1239}%
\special{ip}%
\special{pa 467 1267}%
\special{pa 473 1273}%
\special{ip}%
\special{pa 500 1300}%
\special{pa 506 1306}%
\special{ip}%
\special{pa 533 1333}%
\special{pa 539 1339}%
\special{ip}%
\special{pa 567 1367}%
\special{pa 573 1373}%
\special{ip}%
\special{ip}%
\special{pa 600 1400}%
\special{pa 1400 2200}%
\special{fp}%
\special{pn 8}%
\special{pa 1406 2206}%
\special{pa 1433 2233}%
\special{ip}%
\special{pa 1439 2239}%
\special{pa 1467 2267}%
\special{ip}%
\special{pa 1473 2273}%
\special{pa 1500 2300}%
\special{ip}%
\special{pa 1506 2306}%
\special{pa 1533 2333}%
\special{ip}%
\special{pa 1539 2339}%
\special{pa 1567 2367}%
\special{ip}%
\special{pa 1573 2373}%
\special{pa 1600 2400}%
\special{ip}%
% LINE 2 2 3 0 Black White  
% 4 2370 1080 1170 1080 1170 1080 1170 2400
% 
\special{pn 8}%
\special{pa 2370 1080}%
\special{pa 1170 1080}%
\special{dt 0.045}%
\special{pa 1170 1080}%
\special{pa 1170 2400}%
\special{dt 0.045}%
% STR 2 0 3 0 Black White  
% 4 2200 1300 2200 1400 2 0 1 0
% $A$
\put(22.0000,-14.0000){\makebox(0,0)[lb]{{\colorbox[named]{White}{$A$}}}}%
% STR 2 0 3 0 Black White  
% 4 1400 500 1400 600 2 0 1 0
% $D$
\put(14.0000,-6.0000){\makebox(0,0)[lb]{{\colorbox[named]{White}{$D$}}}}%
% STR 2 0 3 0 Black White  
% 4 600 1300 600 1400 2 0 1 0
% $C$
\put(6.0000,-14.0000){\makebox(0,0)[lb]{{\colorbox[named]{White}{$C$}}}}%
% STR 2 0 3 0 Black White  
% 4 1170 980 1170 1080 2 0 1 0
% $P$
\put(11.7000,-10.8000){\makebox(0,0)[lb]{{\colorbox[named]{White}{$P$}}}}%
% STR 2 0 3 0 Black White  
% 4 1880 980 1880 1080 2 0 1 0
% $Q$
\put(18.8000,-10.8000){\makebox(0,0)[lb]{{\colorbox[named]{White}{$Q$}}}}%
% LINE 3 0 3 0 Black White  
% 24 1840 1400 1400 1840 1900 1400 1400 1900 1960 1400 1400 1960 2020 1400 1400 2020 2080 1400 1400 2080 2140 1400 1400 2140 1780 1400 1400 1780 1720 1400 1400 1720 1660 1400 1400 1660 1600 1400 1400 1600 1540 1400 1400 1540 1480 1400 1400 1480
% 
\special{pn 4}%
\special{pa 1840 1400}%
\special{pa 1400 1840}%
\special{fp}%
\special{pa 1900 1400}%
\special{pa 1400 1900}%
\special{fp}%
\special{pa 1960 1400}%
\special{pa 1400 1960}%
\special{fp}%
\special{pa 2020 1400}%
\special{pa 1400 2020}%
\special{fp}%
\special{pa 2080 1400}%
\special{pa 1400 2080}%
\special{fp}%
\special{pa 2140 1400}%
\special{pa 1400 2140}%
\special{fp}%
\special{pa 1780 1400}%
\special{pa 1400 1780}%
\special{fp}%
\special{pa 1720 1400}%
\special{pa 1400 1720}%
\special{fp}%
\special{pa 1660 1400}%
\special{pa 1400 1660}%
\special{fp}%
\special{pa 1600 1400}%
\special{pa 1400 1600}%
\special{fp}%
\special{pa 1540 1400}%
\special{pa 1400 1540}%
\special{fp}%
\special{pa 1480 1400}%
\special{pa 1400 1480}%
\special{fp}%
% LINE 3 0 3 0 Black White  
% 36 1800 1080 1480 1400 1740 1080 1420 1400 1680 1080 1400 1360 1620 1080 1400 1300 1560 1080 1400 1240 1500 1080 1400 1180 1440 1080 1400 1120 1860 1080 1540 1400 1900 1100 1600 1400 1930 1130 1660 1400 1960 1160 1720 1400 1990 1190 1780 1400 2020 1220 1840 1400 2050 1250 1900 1400 2080 1280 1960 1400 2110 1310 2020 1400 2140 1340 2080 1400 2170 1370 2140 1400
% 
\special{pn 4}%
\special{pa 1800 1080}%
\special{pa 1480 1400}%
\special{fp}%
\special{pa 1740 1080}%
\special{pa 1420 1400}%
\special{fp}%
\special{pa 1680 1080}%
\special{pa 1400 1360}%
\special{fp}%
\special{pa 1620 1080}%
\special{pa 1400 1300}%
\special{fp}%
\special{pa 1560 1080}%
\special{pa 1400 1240}%
\special{fp}%
\special{pa 1500 1080}%
\special{pa 1400 1180}%
\special{fp}%
\special{pa 1440 1080}%
\special{pa 1400 1120}%
\special{fp}%
\special{pa 1860 1080}%
\special{pa 1540 1400}%
\special{fp}%
\special{pa 1900 1100}%
\special{pa 1600 1400}%
\special{fp}%
\special{pa 1930 1130}%
\special{pa 1660 1400}%
\special{fp}%
\special{pa 1960 1160}%
\special{pa 1720 1400}%
\special{fp}%
\special{pa 1990 1190}%
\special{pa 1780 1400}%
\special{fp}%
\special{pa 2020 1220}%
\special{pa 1840 1400}%
\special{fp}%
\special{pa 2050 1250}%
\special{pa 1900 1400}%
\special{fp}%
\special{pa 2080 1280}%
\special{pa 1960 1400}%
\special{fp}%
\special{pa 2110 1310}%
\special{pa 2020 1400}%
\special{fp}%
\special{pa 2140 1340}%
\special{pa 2080 1400}%
\special{fp}%
\special{pa 2170 1370}%
\special{pa 2140 1400}%
\special{fp}%
% LINE 3 0 3 0 Black White  
% 32 1400 1600 1170 1830 1400 1660 1170 1890 1400 1720 1170 1950 1400 1780 1190 1990 1400 1840 1220 2020 1400 1900 1250 2050 1400 1960 1280 2080 1400 2020 1310 2110 1400 2080 1340 2140 1400 2140 1370 2170 1400 1540 1170 1770 1400 1480 1170 1710 1400 1420 1170 1650 1360 1400 1170 1590 1300 1400 1170 1530 1240 1400 1170 1470
% 
\special{pn 4}%
\special{pa 1400 1600}%
\special{pa 1170 1830}%
\special{fp}%
\special{pa 1400 1660}%
\special{pa 1170 1890}%
\special{fp}%
\special{pa 1400 1720}%
\special{pa 1170 1950}%
\special{fp}%
\special{pa 1400 1780}%
\special{pa 1190 1990}%
\special{fp}%
\special{pa 1400 1840}%
\special{pa 1220 2020}%
\special{fp}%
\special{pa 1400 1900}%
\special{pa 1250 2050}%
\special{fp}%
\special{pa 1400 1960}%
\special{pa 1280 2080}%
\special{fp}%
\special{pa 1400 2020}%
\special{pa 1310 2110}%
\special{fp}%
\special{pa 1400 2080}%
\special{pa 1340 2140}%
\special{fp}%
\special{pa 1400 2140}%
\special{pa 1370 2170}%
\special{fp}%
\special{pa 1400 1540}%
\special{pa 1170 1770}%
\special{fp}%
\special{pa 1400 1480}%
\special{pa 1170 1710}%
\special{fp}%
\special{pa 1400 1420}%
\special{pa 1170 1650}%
\special{fp}%
\special{pa 1360 1400}%
\special{pa 1170 1590}%
\special{fp}%
\special{pa 1300 1400}%
\special{pa 1170 1530}%
\special{fp}%
\special{pa 1240 1400}%
\special{pa 1170 1470}%
\special{fp}%
% LINE 3 0 3 0 Black White  
% 18 1400 1180 1180 1400 1400 1120 1170 1350 1380 1080 1170 1290 1320 1080 1170 1230 1260 1080 1170 1170 1200 1080 1170 1110 1400 1240 1240 1400 1400 1300 1300 1400 1400 1360 1360 1400
% 
\special{pn 4}%
\special{pa 1400 1180}%
\special{pa 1180 1400}%
\special{fp}%
\special{pa 1400 1120}%
\special{pa 1170 1350}%
\special{fp}%
\special{pa 1380 1080}%
\special{pa 1170 1290}%
\special{fp}%
\special{pa 1320 1080}%
\special{pa 1170 1230}%
\special{fp}%
\special{pa 1260 1080}%
\special{pa 1170 1170}%
\special{fp}%
\special{pa 1200 1080}%
\special{pa 1170 1110}%
\special{fp}%
\special{pa 1400 1240}%
\special{pa 1240 1400}%
\special{fp}%
\special{pa 1400 1300}%
\special{pa 1300 1400}%
\special{fp}%
\special{pa 1400 1360}%
\special{pa 1360 1400}%
\special{fp}%
% STR 2 0 3 0 Black White  
% 4 1170 1880 1170 1980 2 0 1 0
% $R$
\put(11.7000,-19.8000){\makebox(0,0)[lb]{{\colorbox[named]{White}{$R$}}}}%
% STR 2 0 3 0 Black White  
% 4 1400 2100 1400 2200 2 0 1 0
% $B$
\put(14.0000,-22.0000){\makebox(0,0)[lb]{{\colorbox[named]{White}{$B$}}}}%
\end{picture}}%
}
% \end{center}

$0\le x$の時，明らかに折れ線は正方形の面積を二等分しないから
\begin{align}
  t-1\le x<0 \label{eq:1}
\end{align}
である．この時，折れ線と$BC$の交点は$R(x,-x-1)$である．$PR$と$x$軸の交点$T$とすれば，題意の面積についての条件から，
\begin{align}
  &\triangle CRT=\square AQPT \\
  &\dfrac{1}{2}(1+x)^2=\dfrac{1}{2}t\left((1-t-x)+(1-x)\right) \\
  &=\dfrac{1}{2}t(2-2x-t) \\
  &x^2+2x+1=2t-2xt-t^2 \\
  &x^2+t^2+2xt+2x-2t+1=0 \\
  &x^2+2(t+1)x+(t-1)^2=0 \\
  &x=-(t+1)\pm\sqrt{(t+1)^2-(t-1)^2} \\
  &x=-(t+1)\pm2\sqrt{t} \\
\end{align}
$P$は$(0,1)$を通るので，複号正をとって
\begin{align}
  x=-t-1+2\sqrt{t}=-(\sqrt{t}-1)^2
\end{align}
である．よってグラフの概形は下図である．

% TODO: 手書きpdfの作図
% \begin{center}
%   \scalebox{1}{%WinTpicVersion4.32a
{\unitlength 0.1in%
\begin{picture}(20.0000,10.4000)(4.0000,-14.4000)%
% STR 2 0 3 0 Black White  
% 4 1390 1397 1390 1410 4 400 0 0
% O
\put(13.9000,-14.1000){\makebox(0,0)[rt]{O}}%
% STR 2 0 3 0 Black White  
% 4 1360 387 1360 400 4 400 0 0
% $y$
\put(13.6000,-4.0000){\makebox(0,0)[rt]{$y$}}%
% STR 2 0 3 0 Black White  
% 4 2400 1427 2400 1440 4 400 0 0
% $x$
\put(24.0000,-14.4000){\makebox(0,0)[rt]{$x$}}%
% VECTOR 2 0 3 0 Black White  
% 2 1400 1400 1400 400
% 
\special{pn 8}%
\special{pa 1400 1400}%
\special{pa 1400 400}%
\special{fp}%
\special{sh 1}%
\special{pa 1400 400}%
\special{pa 1380 467}%
\special{pa 1400 453}%
\special{pa 1420 467}%
\special{pa 1400 400}%
\special{fp}%
% VECTOR 2 0 3 0 Black White  
% 2 400 1400 2400 1400
% 
\special{pn 8}%
\special{pa 400 1400}%
\special{pa 2400 1400}%
\special{fp}%
\special{sh 1}%
\special{pa 2400 1400}%
\special{pa 2333 1380}%
\special{pa 2347 1400}%
\special{pa 2333 1420}%
\special{pa 2400 1400}%
\special{fp}%
% FUNC 2 0 3 0 Black White  
% 9 400 400 2400 1400 1400 1400 2200 1400 1400 600 1400 400 2200 1400 0 4 0 0
% -x+1
\special{pn 8}%
\special{pn 8}%
\special{pa 1200 400}%
\special{pa 1206 406}%
\special{ip}%
\special{pa 1233 433}%
\special{pa 1239 439}%
\special{ip}%
\special{pa 1267 467}%
\special{pa 1273 473}%
\special{ip}%
\special{pa 1300 500}%
\special{pa 1306 506}%
\special{ip}%
\special{pa 1333 533}%
\special{pa 1339 539}%
\special{ip}%
\special{pa 1367 567}%
\special{pa 1373 573}%
\special{ip}%
\special{ip}%
\special{pa 1400 600}%
\special{pa 2200 1400}%
\special{fp}%
\special{pn 8}%
% FUNC 2 0 3 0 Black White  
% 10 400 400 2400 1400 1400 1400 2200 1400 1400 600 600 400 1400 1400 0 4 0 0 0 0
% x+1
\special{pn 8}%
\special{pn 8}%
\special{pa 600 1400}%
\special{pa 600 1400}%
\special{ip}%
\special{pa 600 1400}%
\special{pa 1400 600}%
\special{fp}%
\special{pn 8}%
\special{pa 1406 594}%
\special{pa 1433 567}%
\special{ip}%
\special{pa 1439 561}%
\special{pa 1467 533}%
\special{ip}%
\special{pa 1473 527}%
\special{pa 1500 500}%
\special{ip}%
\special{pa 1506 494}%
\special{pa 1533 467}%
\special{ip}%
\special{pa 1539 461}%
\special{pa 1567 433}%
\special{ip}%
\special{pa 1573 427}%
\special{pa 1600 400}%
\special{ip}%
% FUNC 2 0 3 0 Black White  
% 10 400 400 2400 1400 1400 1400 2200 1400 1400 600 400 600 2400 1400 0 4 0 1 0 0
% -y-1+2sqrt(y)
\special{pn 8}%
\special{pn 8}%
\special{pa 1389 400}%
\special{pa 1389 405}%
\special{pa 1390 409}%
\special{ip}%
\special{pa 1394 450}%
\special{pa 1394 450}%
\special{pa 1394 459}%
\special{ip}%
\special{pa 1397 500}%
\special{pa 1397 505}%
\special{pa 1398 509}%
\special{ip}%
\special{pa 1399 550}%
\special{pa 1399 555}%
\special{pa 1400 559}%
\special{ip}%
\special{pa 1400 600}%
\special{pa 1400 635}%
\special{pa 1399 640}%
\special{pa 1399 665}%
\special{pa 1398 670}%
\special{pa 1398 685}%
\special{pa 1397 690}%
\special{pa 1397 700}%
\special{pa 1396 705}%
\special{pa 1396 715}%
\special{pa 1395 720}%
\special{pa 1395 725}%
\special{pa 1394 730}%
\special{pa 1394 735}%
\special{pa 1393 740}%
\special{pa 1393 745}%
\special{pa 1392 750}%
\special{pa 1392 755}%
\special{pa 1390 765}%
\special{pa 1390 770}%
\special{pa 1389 775}%
\special{pa 1389 780}%
\special{pa 1386 795}%
\special{pa 1386 800}%
\special{pa 1381 825}%
\special{pa 1381 830}%
\special{pa 1375 860}%
\special{pa 1373 865}%
\special{pa 1369 885}%
\special{pa 1367 890}%
\special{pa 1364 905}%
\special{pa 1362 910}%
\special{pa 1361 915}%
\special{pa 1359 920}%
\special{pa 1358 925}%
\special{pa 1356 930}%
\special{pa 1355 935}%
\special{pa 1353 940}%
\special{pa 1352 945}%
\special{pa 1348 955}%
\special{pa 1347 960}%
\special{pa 1323 1020}%
\special{pa 1320 1025}%
\special{pa 1316 1035}%
\special{pa 1313 1040}%
\special{pa 1311 1045}%
\special{pa 1308 1050}%
\special{pa 1306 1055}%
\special{pa 1300 1065}%
\special{pa 1298 1070}%
\special{pa 1277 1105}%
\special{pa 1273 1110}%
\special{pa 1267 1120}%
\special{pa 1263 1125}%
\special{pa 1260 1130}%
\special{pa 1224 1175}%
\special{pa 1214 1185}%
\special{pa 1210 1190}%
\special{pa 1190 1210}%
\special{pa 1184 1215}%
\special{pa 1179 1220}%
\special{pa 1173 1225}%
\special{pa 1168 1230}%
\special{pa 1156 1240}%
\special{pa 1149 1245}%
\special{pa 1143 1250}%
\special{pa 1115 1270}%
\special{pa 1107 1275}%
\special{pa 1100 1280}%
\special{pa 1092 1285}%
\special{pa 1083 1290}%
\special{pa 1075 1295}%
\special{pa 1066 1300}%
\special{pa 1056 1305}%
\special{pa 1047 1310}%
\special{pa 1037 1315}%
\special{pa 1015 1325}%
\special{pa 991 1335}%
\special{pa 965 1345}%
\special{pa 950 1350}%
\special{pa 918 1360}%
\special{pa 900 1365}%
\special{pa 880 1370}%
\special{pa 858 1375}%
\special{pa 833 1380}%
\special{pa 804 1385}%
\special{pa 769 1390}%
\special{pa 721 1395}%
\special{pa 600 1400}%
\special{fp}%
% LINE 3 0 3 0 Black White  
% 52 1900 1100 1600 1400 1870 1070 1540 1400 1840 1040 1480 1400 1810 1010 1420 1400 1780 980 1400 1360 1750 950 1400 1300 1720 920 1400 1240 1690 890 1400 1180 1660 860 1400 1120 1630 830 1400 1060 1600 800 1400 1000 1570 770 1400 940 1540 740 1400 880 1510 710 1400 820 1480 680 1400 760 1450 650 1400 700 1420 620 1400 640 1930 1130 1660 1400 1960 1160 1720 1400 1990 1190 1780 1400 2020 1220 1840 1400 2050 1250 1900 1400 2080 1280 1960 1400 2110 1310 2020 1400 2140 1340 2080 1400 2170 1370 2140 1400
% 
\special{pn 4}%
\special{pa 1900 1100}%
\special{pa 1600 1400}%
\special{fp}%
\special{pa 1870 1070}%
\special{pa 1540 1400}%
\special{fp}%
\special{pa 1840 1040}%
\special{pa 1480 1400}%
\special{fp}%
\special{pa 1810 1010}%
\special{pa 1420 1400}%
\special{fp}%
\special{pa 1780 980}%
\special{pa 1400 1360}%
\special{fp}%
\special{pa 1750 950}%
\special{pa 1400 1300}%
\special{fp}%
\special{pa 1720 920}%
\special{pa 1400 1240}%
\special{fp}%
\special{pa 1690 890}%
\special{pa 1400 1180}%
\special{fp}%
\special{pa 1660 860}%
\special{pa 1400 1120}%
\special{fp}%
\special{pa 1630 830}%
\special{pa 1400 1060}%
\special{fp}%
\special{pa 1600 800}%
\special{pa 1400 1000}%
\special{fp}%
\special{pa 1570 770}%
\special{pa 1400 940}%
\special{fp}%
\special{pa 1540 740}%
\special{pa 1400 880}%
\special{fp}%
\special{pa 1510 710}%
\special{pa 1400 820}%
\special{fp}%
\special{pa 1480 680}%
\special{pa 1400 760}%
\special{fp}%
\special{pa 1450 650}%
\special{pa 1400 700}%
\special{fp}%
\special{pa 1420 620}%
\special{pa 1400 640}%
\special{fp}%
\special{pa 1930 1130}%
\special{pa 1660 1400}%
\special{fp}%
\special{pa 1960 1160}%
\special{pa 1720 1400}%
\special{fp}%
\special{pa 1990 1190}%
\special{pa 1780 1400}%
\special{fp}%
\special{pa 2020 1220}%
\special{pa 1840 1400}%
\special{fp}%
\special{pa 2050 1250}%
\special{pa 1900 1400}%
\special{fp}%
\special{pa 2080 1280}%
\special{pa 1960 1400}%
\special{fp}%
\special{pa 2110 1310}%
\special{pa 2020 1400}%
\special{fp}%
\special{pa 2140 1340}%
\special{pa 2080 1400}%
\special{fp}%
\special{pa 2170 1370}%
\special{pa 2140 1400}%
\special{fp}%
% LINE 3 0 3 0 Black White  
% 26 1400 1180 1180 1400 1400 1120 1120 1400 1400 1060 1060 1400 1190 1210 1000 1400 1020 1320 940 1400 920 1360 880 1400 1400 1000 1210 1190 1400 940 1320 1020 1400 880 1350 930 1400 820 1370 850 1400 1240 1240 1400 1400 1300 1300 1400 1400 1360 1360 1400
% 
\special{pn 4}%
\special{pa 1400 1180}%
\special{pa 1180 1400}%
\special{fp}%
\special{pa 1400 1120}%
\special{pa 1120 1400}%
\special{fp}%
\special{pa 1400 1060}%
\special{pa 1060 1400}%
\special{fp}%
\special{pa 1190 1210}%
\special{pa 1000 1400}%
\special{fp}%
\special{pa 1020 1320}%
\special{pa 940 1400}%
\special{fp}%
\special{pa 920 1360}%
\special{pa 880 1400}%
\special{fp}%
\special{pa 1400 1000}%
\special{pa 1210 1190}%
\special{fp}%
\special{pa 1400 940}%
\special{pa 1320 1020}%
\special{fp}%
\special{pa 1400 880}%
\special{pa 1350 930}%
\special{fp}%
\special{pa 1400 820}%
\special{pa 1370 850}%
\special{fp}%
\special{pa 1400 1240}%
\special{pa 1240 1400}%
\special{fp}%
\special{pa 1400 1300}%
\special{pa 1300 1400}%
\special{fp}%
\special{pa 1400 1360}%
\special{pa 1360 1400}%
\special{fp}%
\end{picture}}%
}
% \end{center}

求める面積$S$は
\begin{align}
  S&=\int_0^1(1-t+(\sqrt{t}-1)^2)dt \\
  &=\int_0^12(1-\sqrt{t})dt \\
  &=2\left[t-\dfrac{2}{3}t^{3/2}\right]_0^1 \\
  &=\dfrac{2}{3}
\end{align}
である．一辺$a$とすると，これの$a^2/2$倍ゆえ
\begin{align}
  S=\dfrac{1}{3}a^2
\end{align}
である．$\cdots$(答)
\newpage
\end{document}
